%% kapa: exam=1 round=35 year=2024 period=2 form=A subject=appraisal_law subject_ko=감정평가관계법규 seq=30 number=30
%% source: Q-Net 감정평가사 기출문제 (cst00309) artlSeq=5213229 file=raw/1008/5213229/2교시 문제지 원본.pdf
%% conversion: kordoc 4.24.51 -> tools/gen1.py (auto); review=proofread-sonnet; flags=none
\begin{question}{30}
공간정보의 구축 및 관리 등에 관한 법령상 지목에 관한 설명으로 옳은 것을 모두 고른 것은?
\begin{framedtext}
ㄱ. 지목의 설정은 필지마다 하나의 지목을 설정하는 방법으로 한다. \par
ㄴ. 송유시설의 부지는 지목을 잡종지로 한다. \par
ㄷ. 건축물의 용도가 변경된 경우는 지목변경을 신청할 수 없다. \par
ㄹ. 지적소관청은 지목변경을 하려면 시ㆍ도지사의 승인을 받아야 한다.
\end{framedtext}
\begin{choices}
  \item ㄱ, ㄴ
  \item ㄱ, ㄷ
  \item ㄷ, ㄹ
  \item ㄱ, ㄴ, ㄹ
  \item ㄴ, ㄷ, ㄹ
\end{choices}
\end{question}
\answer{1}
